\section{Study Design}
\label{sec:design}

%% Revised 2026-08-30 against the GPT-5.6-Sol review (plans/2026-08-30-saner-review-response-plan.md)
%% and tightened for the page budget. Numbers are macros only.
%% 2026-09-13 shortening pass (response plan §13): the terminology paragraph is the shared
%% latex/shared/terms.tex (also used by the TR; captions cite it as \S\ref{sec:design}); M5/M7
%% send the full inferential machinery, the metric definitions and the entity-ownership
%% tie-breaks to the TR; caveats stated here are not restated in §V/§VI.
%% 2026-09-13 readability pass: long sentences split, §IV.D split into two paragraphs; no content change.

The study is built so the two-regime thesis of \S\ref{sec:intro} can fail.
\emph{Recoverability} (RQ1--RQ3) asks how much of a system's module structure is recoverable from
build and dependency evidence alone, and where that breaks down. \emph{Expression and governance}
(RQ4--RQ7) asks whether, where recovery fails, a human-in-the-loop pipeline keeps the expressed
architecture stable, checkable, and affordable in a way autonomous SAR does not.
\textbf{RQ1:} how much does build evidence alone recover, can that be diagnosed without a
reference, and where is human input needed? The autonomous rows are the \emph{no-curation floor}
and the \emph{auto-propose draft} (defined below). \textbf{RQ2:} given the same dependency input,
do the autonomous rows match classical SAR and a controlled LLM baseline, and how far does
curation move the score? \textbf{RQ3:} which design choices contribute, at the operating
configuration and at the extremes? \textbf{RQ4:} are outputs identical across reruns, and does
the decomposition survive code change? \textbf{RQ5:} does drift detection report seeded and real
change at an acceptable false-finding cost, and fail honestly when evidence goes missing?
\textbf{RQ6:} how do runtime and memory scale against the pipeline's own budget? \textbf{RQ7:}
where recovery fails, how close does a reference-blind curator get to the reference, at what
effort, and how does that compare with the agreement between two reference authors? The plan's
RQ identifiers are mapped in the artifact.

%% The shared terminology paragraph (was the boxed Table I until 2026-09-13; folded into prose
%% for the page budget, response plan section 13 status). The legend the table captions used to
%% carry lives here once; captions reference \S\ref{sec:design} (the enclosing section: IV in the
%% paper, 1 in the TR, which labels its scope section sec:design too) instead of restating it.
%% SHARED between main.tex and tr.tex: it must not reference a label that exists in only one of
%% the two documents. Hand-owned: definitions, not data.
\paragraph{Terminology.}
Three \emph{rows} recur in every table: the \emph{floor} (autonomous: every scored entity in its
owning build target, no curation), the \emph{propose} draft (autonomous: the tool's heuristic rule
draft accepted verbatim), and the \emph{curated} model (the committed hand rules applied). A
\emph{detection} number scores autonomous recovery, an \emph{expression} number what a human
expressed under evidence constraints; the two are never compared inferentially, and only detection
numbers support a recovery claim. Curation provenance (Cur.): \emph{sighted} = reference in view,
so the score is an oracle-assisted upper bound, never evidence against an autonomous technique;
\emph{assisted} = LLM-assisted from names and directories, the reference mapping not consulted;
\emph{blind} = reference-blind curator (RQ7); \emph{none} = curated model identical to the floor.
Reference grades (Ref) and the \emph{all-common} projection are defined in
\S\ref{sec:corpus} and \S\ref{sec:baselines}.

Curation sets the very partition the metrics score, so a curated score compares assisted with
unassisted recovery and says nothing about autonomous recovery. Every detection claim rests on
the two autonomous rows; the non-circular curated numbers come from the reference-blind
curations (RQ7; Squidex's only curated model). Three of the four C\#\ curations
were authored with the reference in view, and P-1's with the diagram its reference was
reconciled from (Table~\ref{tab:corpus}).

\subsection{Subject corpus and references}
\label{sec:corpus}

\textbf{Tier-A} (Nine systems: five C/C++, four C\#) carries
a reference architecture and drives RQ1--RQ3 and RQ7 (Table~\ref{tab:corpus}); \textbf{Tier-B}
(\rqSevenSystems{} systems, containing Tier-A) needs none and drives RQ4--RQ6. The anonymized
industrial system, \textbf{P-1} in the tables, sits outside both tiers: its metrics corroborate
the open-source results but are not replicable. Tier-B spans three orders of magnitude in KLOC and
four build idioms (CMake, .NET SDK, MSBuild-C++, automake or hand-written \texttt{Makefile.in}),
which co-vary with language here (\S\ref{sec:threats}). Every system is pinned to a commit;
nopCommerce's source (license) and libxml2's reference (third-party) stay out of our package,
their measurements in. \textbf{GT-pub} adopts a published mapping verbatim. \textbf{GT-doc} is transcribed from the project's own documentation by two
people independently, then reconciled. For \textbf{GT-exp}, two architects map entities to
components blind to the pipeline output; we report their agreement (\S\ref{sec:metrics}), resolve
by discussion, and flag $\kappa<0.6$ as lower-confidence. abseil is the exception: a single
rater applied its documented one-library-per-directory layout, so no agreement is reported. No
system's curator authors its reference. Table~\ref{tab:corpus} records each curation's provenance (attested; overlay and
reference commit dates in the artifact).

\subsection{Baselines, their inputs, and fairness controls}
\label{sec:baselines}

We compare against \rqTwoBaselineCount{} classical SAR and na\"ive baselines plus a controlled
LLM-recovery baseline. The semantic and agentic lines of \S\ref{sec:background} are discussed, not
run: SARIF's only release is a closed Linux binary without C\#\ support, FCA is an unlicensed
MATLAB script, and the four most recent comparative evaluations themselves take this classical
set as their baselines~\cite{zhang2023sarif,zhao2026semarc,ding2026ssar,zhang2026semref}. The
classical set is ACDC~\cite{acdc}, ARC~\cite{arc}, WCA~\cite{wca}, and LIMBO~\cite{limbo}
(\S\ref{sec:background}), all run through the
\textsc{Arcade} workbench~\cite{arcade}, which also supplies the MoJoFM and a2a implementations
(c2c is not reported). Two na\"ive lower
bounds, directory-tree clustering and connected components, complete the set. The
\textbf{LLM-recovery baseline} stands for the LLM-centric line~\cite{archagent}: a reasoning-tier
model (\texttt{gpt-5.6-sol}, pinned) receives the exported graph, one line of repository
metadata (the system's name and language), and the same $k$ as ARC under one prompt asking for a JSON entity-to-cluster map, at
default sampling and a 180\,k-token input budget; five reruns on \rqTwoLlmSystemsRan{} systems at
most \$\rqTwoLlmCostMaxUsd{} per system, mean and spread reported, prompt and replies in the
artifact (\trInputs{}).

\paragraph{What each technique receives.} One dependency graph per system, exported once as RSF at
the scored granularity, feeds every technique: at project granularity (C\#) every relationship
kind between first-party projects as one unweighted edge set; at file granularity (C++) the
\texttt{\#include} graph, declared link edges never projected onto files. ACDC and
\texttt{cc} see the edges only; \texttt{dir} sees entity paths; WCA and LIMBO also receive $k$,
the reference cluster count, as their stopping point; ARC receives $k$ and each entity's source
text; the LLM receives the graph, the metadata line, and $k$. $k$ is an oracle hint no pipeline
row gets. The pipeline's autonomous rows instead know each entity's \emph{owning build target},
the signal under test, which no baseline receives in any form. A file belongs to the target whose
source list names it, else to the target owning the longest directory prefix; files no target
claims stay unclustered and count against the floor (\trInputs{}). RQ2 therefore compares the
whole pipeline with clustering the same edges, not two algorithms on identical information. No
baseline receives convention tags, documentation, or mapping rules; the curated row does, hence
its expression label.

\paragraph{Fairness controls, fixed before scoring.} One granularity per comparison (file for
C++, project for C\#), the pipeline's per-entity cluster induced via
\texttt{entity}$\to$\texttt{target}$\to$\texttt{container}; the same reference and metric
implementation for every technique; default parameters; and an \emph{all-common} entity
projection on which every number in Tables~\ref{tab:rq1} and~\ref{tab:rq2} is scored, floor and
propose included. ARC abstains on low-text files; it is scored on its own covered subset and kept
outside every paired test. A technique that does not complete is reported as missing. \emph{Lead}
is defined, not asserted: a row leads a system when it strictly exceeds every all-common classical
baseline (ACDC, WCA, LIMBO, \texttt{dir}, \texttt{cc}) on that metric; wins-or-ties and wins
against all displayed methods are reported beside it. The \emph{recovery bands} are post hoc:
recovered at MoJoFM $\ge\rqTwoRecoveredMin{}$, the score the reference-free diagnostic of
\S\ref{sec:metrics} is checked against, failed below \rqTwoFailedBelow{}, partial between.

\subsection{Ablations and blind curation}
\label{sec:ablation}

Eight axes toggle independently from the full configuration, paired per system on Tier-A: A1 caps
the evidence ladder at L0, L1, or L2; A2 walks the curation rungs (floor, single heuristics, the
auto-propose draft, hand rules); A3 sweeps the weight threshold; A4 toggles visibility
down-weighting; A5 switches off the never-drop-declared rule at the operating threshold and above;
A6 reads the RQ2 LLM-recovery runs as the LLM-does-structure arm; A7 toggles LLM naming; A8 measures stability under code
change. A frozen \texttt{mapping-rules.yaml} per system holds curation constant.

The blind-curation experiment (RQ7) yields the curation-vs-reference numbers that are not
circular. Its protocol, roles, and stopping rule were written into the plan before any session
and are unchanged since (dated commit history and signed attestations in the artifact; no external
registry). It ran on two C\#\ systems with four researchers in disjoint roles: eShop, whose
documented reference two raters re-derived independently as the leakage check against the sighted
curation, and Squidex, whose reference the raters built. Reference authors R$_1$ and R$_2$ work
from code and project docs. A curator C, who is neither, authors the mapping rules blind to the
reference, to R's mappings, and to any score, using the repository, docs, \cli~\texttt{propose}
draft, build facts, and GUI. C stops when the model is review-complete by their own judgment;
the harness logs the edit trace and wall-clock. Each system yields detection rows (the floor, the
best all-common technique), the expression row (the blind curator against the reference), and the
R$_1$-vs-R$_2$ agreement, an \emph{inter-rater baseline} that bounds nothing about a third
partition.

\subsection{Metrics and analysis}
\label{sec:metrics}

We lead with a2a~\cite{le2015a2a} and report MoJoFM~\cite{mojofm} second. MoJoFM rewards matching
the reference's cluster count, and a2a charges for the sprawl that produces, which is why the
na\"ive \texttt{dir} baseline beats the pipeline on MoJoFM on directory-aligned C++ but loses on
a2a. At a reviewer's request we add, post hoc and descriptively, the adjusted Rand index (ARI,
$\times$100), chance-corrected and charging over-splitting in full~\cite{zhang2023sarif}, and the
\emph{nested share}, the fraction of a partition's clusters lying inside one reference cluster,
which tells over-splitting from misgrouping. The inferential tests stay on MoJoFM and a2a.
\emph{Container-edge F1} scores the lifted container edges against the reference's own edge set
after a greedy label alignment; it is curation-constrained, not curation-proof, because the
grouping fixes which edges become intra-container. Inter-rater agreement is reported label-free
(ARI, normalized mutual information, and Cohen's $\kappa$ after Hungarian label matching), because
a raw label-match $\kappa$ measures wording. The raw $\kappa$ stays beside them because the plan's
flag rule (raw $\kappa<0.6$ marks a low-confidence reference) is applied as written.

To diagnose ``the build graph carries the architecture'' \emph{without} a reference we compute,
from the extracted facts alone, the floor's \emph{unit ratio} (build units per scored entity;
1.0 = every entity its own unit), its \emph{largest-unit share}, and its TurboMQ over the exported
graph. The rule over them was selected on the same systems it separates, so beside Spearman's
$\rho$ with a bootstrap interval we report the threshold interval over which it keeps its
agreement and a leave-one-system-out re-selection (\trDiagnostic{}). Governance is measured by
byte-identical regeneration and survivor-partition agreement at 20\,\% turnover (RQ4). Drift
detection (RQ5; \trDrift{}) is measured by precision and recall against seeded build-file
mutations run through the whole pipeline, per mutation and per finding; by the suppression rate
under induced evidence degradation; by the outcome of every seeded finding when a change and
evidence loss coincide; and by the findings between consecutive real releases under rules held
constant across each pair, each labelled by two raters. Performance (RQ6) is per-stage
wall-clock and peak RSS.

Scripts and thresholds were committed before results were unblinded. Across techniques (RQ2) the
inferential tests admit autonomous rows only: a Friedman omnibus on MoJoFM and a2a over the
all-common techniques with the floor (and, separately, the propose draft) in place of the curated
row, then Nemenyi and Holm-corrected pairwise Wilcoxon signed-rank tests of that row against each
technique (\trStats{}). The curated row enters no inferential test; its omnibus is descriptive. Paired on/off comparisons (RQ3) use Wilcoxon
signed-rank with the matched-pairs rank-biserial $r$. Cliff's~$\delta$ is reserved for the one
between-group contrast (C++ vs.\ C\#\ floor); detection rates get Wilson intervals, the stopping
rules exact McNemar, and the diagnostic Spearman's $\rho$. With $N\le 5$ paired systems the
two-sided exact signed-rank $p$ cannot fall below $0.0625$, and at $N=6$ not below $0.0312$,
however large the effect, so effect sizes and per-system deltas carry the argument. Every statistic is a generated macro.
